% Source-derived chapter 05: Topology of Hitchin fibrations
% Original source: hitchin-fibrations-abelian-surfaces-pw
\chapter{Topology of Hitchin fibrations}
\label{hitchin-fibrations-abelian-surfaces-pw-chapter-05}
\source{hitchin-fibrations-abelian-surfaces-pw}

\begin{remark}[Source section]
\label{hitchin-fibrations-abelian-surfaces-pw-chapter-05-source}
\source{hitchin-fibrations-abelian-surfaces-pw}
\source{hitchin-fibrations-abelian-surfaces-pw}
This chapter reproduces the corresponding section of the retrieved
LaTeX source, including its definitions, statements, examples, and
proofs. It has not yet been translated into Lean.
\notready
\end{remark}

% The source section title was: Topology of Hitchin fibrations
\label{se4}
\source{hitchin-fibrations-abelian-surfaces-pw}

\subsection{Overview} Throughout the section, we assume $C$ is a nonsingular projective integral curve of genus $g$ embedded into an abelian surface $A$. We study {a kind of} specialization morphism
\begin{equation}\label{sp}
\source{hitchin-fibrations-abelian-surfaces-pw}
    \mathrm{sp}^!: H^*(\CM_{\beta, A}, \BQ) \to H^*(\CM_{\mathrm{Dol}}, \BQ)
\end{equation}
where $\CM_{\mathrm{Dol}}$ is the moduli of stable Higgs bundles of rank $r$ and Euler characteristic $\chi$, and 
\[
\beta = r[C] \in H_2(A, \BZ).
\]
Then we deduce Theorems \ref{genus_2}, \ref{P=W_even}, and \ref{P=W_odd} from Theorem \ref{taut_abelian} via the properties of the  morphism (\ref{sp}) which we establish hereafter.

\subsection{Deformation to the normal cone}\label{Section 4.2}
\source{hitchin-fibrations-abelian-surfaces-pw}

The moduli space $\CM_{\mathrm{Dol}}$ of stable Higgs bundles with rank $r$ and Euler characteristic $\chi$ can be realized as the moduli space of $1$-dimensional Gieseker-stable sheaves $\CF$ on the cotangent bundle  surface $T^*C$ with (cf. see \cite{BNR})
\[
 [\mathrm{supp}(\CF)] = r[C] \in H_2(T^*C, \BZ), \quad \chi(\CF) = \chi.
\]
In the following, we describe $\CM_{\mathrm{Dol}}$ as the ``limit'' of a trivial family of $\CM_{\beta,A}$. A similar construction using $K3$ surfaces was considered in \cite{DLL}.

Assume $T = \BP^1$ and $T^\circ = \BP^1 \setminus 0 = \BC$. Let 
\begin{equation}\label{normal_cone}
\source{hitchin-fibrations-abelian-surfaces-pw}
 \CS = \mathrm{Bl}_{C\times 0}(A \times T) \setminus (A\times 0) \to T
\end{equation}
be the total space of the deformation to the normal cone associated with the embedding $j_C:C\hookrightarrow A$. The central fiber of (\ref{normal_cone}) is
\[
\CS_0 = T^*C \to 0 \in T,
\]
and the restriction over $T^\circ$ is a trivial fibration
\[
A \times T^\circ \to  T^\circ \subset T.
\]


We associate to $\beta$ a family of homology classes 
\begin{equation}\label{beta_t}
\source{hitchin-fibrations-abelian-surfaces-pw}
\beta_t = r[C] \in H_2(\CS_t, \BZ).
\end{equation}
Let $\CM \rightarrow T$ be the (coarse) relative moduli space which parametrizes, for each $t \in T$, pure one-dimensional Gieseker-stable sheaves $\CF$ on $\CS_t$ with $\chi(\CF) =\chi$ and such that the support of $\CF$ is a proper subscheme in the class $\beta_t$. Similarly, let $\CB \rightarrow T$ denote the component of the relative Hilbert scheme which parametrizes Cartier divisors in $\CS_t$ with proper support in the class $\beta_t$.  Following \cite[Section 5.3]{GIT}, we have a proper morphism 
\begin{equation*}
\begin{tikzcd}
\CM \ar [rr,"h_T"]\ar[rd]& & \CB\ar[ld]\\
 &T&
\end{tikzcd}
\end{equation*}  
over the base $T$, which, on the level of points, sends a sheaf $\CF$ on $\CS_t$ to its Fitting support, viewed as an element in $\CB_t$.
Clearly both $\CM$ and $\CB$ become trivial families when restricted over $T^\circ$,
\begin{equation*}
\begin{tikzcd}
\CM_{\beta,A} \times T^\circ \ar [rr,"h_{T^\circ}"]\ar[rd]& & B \times T^\circ \ar[ld]\\
 &T^\circ&.
\end{tikzcd}
\end{equation*}  
The fibers over $0 \in T$ recover the moduli space $\CM_{\mathrm{Dol}}$, the Hitchin base $\Lambda$, and the Hitchin fibration (\ref{Hitchin_fib}).



The following lemma seems standard; we include a brief proof  since we are not aware of an adequate reference.
\begin{lem}
\source{hitchin-fibrations-abelian-surfaces-pw}
\notready\label{lem4.1}
\source{hitchin-fibrations-abelian-surfaces-pw}
\begin{enumerate}

\item[(i)] Both $\CM$ and $\CB$ are irreducible varieties and smooth over $T$.
\item[(ii)] There exists a universal class 
\[
[\CF_T] \in K_{\mathrm{top}}(\CM\times_{T} \CS)
\]
whose restriction to each fiber gives a universal class for $\CM_t \times \CS_t$.
\end{enumerate}

\end{lem}
\begin{proof}

For part (i), we first argue for $\CB$. For notational convenience, we define
\[
\tilde{g} = \mathrm{dim}(\CB_t) = r^2(g-1)+1, \quad \forall~t \in T.
\]
We observe that the closure $\overline\CB^{\circ}$ in $\CB$ of 
\[
\CB^{\circ} = B \times T^\circ
\]
is irreducible of dimension $\tilde{g}+1$, so that the intersection 
\[
\overline{\CB^{\circ}}\cap \CB_0 \subset \CB_0
\]
must have at least dimension $\tilde{g}$ at every closed point.  This intersection is non-empty since it contains the divisor $r[C]$ which clearly deforms to the generic fiber. Since $\CB_0$ is irreducible of dimension $\tilde{g}$, we have $\CB_0 \subset \overline{\CB^{\circ}}$. Therefore $\CB$ is irreducible. This implies the flatness of $\CB$ over the nonsingular curve $T$. Furthermore, since its fibers are all nonsingular, the flatness further implies that the morphism $\CB \to T$ is smooth.

The same argument applies for $\CM$. The only thing to check is that there exists a point in the central fiber $\CM_0$ which deforms to the generic fiber.  By the smoothness of $\CB/T$,  we can choose a nonsingular curve $Z_0 \subset \CS_0$ which deforms to $Z_t \subset \CS_t$.  Any line bundle on $Z_0$ with Euler characteristic $\chi$ can be deformed as well and this gives the desired point in $\CM_0$. 


For part (ii), we follow the same argument from \cite[Section 3.1]{Markman5}. We denote by $\mathfrak{M}$ the relative moduli stack of stable sheaves on $\CS$, which is a $\mathbb{G}_m$-gerbe over $\CM$. It suffices to show that this gerbe is \emph{topologically} trivializable, so that we can pull back the topological $K$-theory class of the universal family on $\mathfrak{M}\times_T\CS$ to $\CM\times_T \CS$ via a section.  To construct a trivialization, we take the nonsingular 3-fold $\CS'= \mathrm{Bl}_{C\times 0}(A \times T)$ which contains $\CS$ as an open subset,
\[
\CS \subset \CS'.
\] 
Let $\CM_{\CS'}$ be the (coarse) moduli space of stable sheaves on $\CS'$ with Chern character given by $w =\mathrm{ch}({i_t}_\ast \CF) \in H^*(\CS', \BQ)$, where 
\[
i_t: \CS_t \xrightarrow{\simeq} A\times t \subset \CS', \quad t\neq 0
\]
is the inclusion of a nonsingular fiber, and $\CF \in \CM$ is a coherent sheaf supported on $\CS_t$. The moduli space $\CM_{\CS'}$ contains $\CM$ as an open subvariety, therefore it is enough to show that the corresponding $\mathbb{G}_m$-gerbe $\mathfrak{M}_{\CS'} \rightarrow \CM_{\CS'}$ is topologically trivial. As in \cite[Section 3.1]{Markman5}, a trivialization can be constructed with a topological K-theory class in $\CS'$ whose pairing with $w$ is $1$. For this, we can use the pullback to $\CS'$ of a class on $A$, whose pairing with $(0 ,r[C], \chi)$ is $1$, via the natural projection
\[
\CS' = \mathrm{Bl}_{C\times 0}(A\times T) \to A. \qedhere
\]
\end{proof}




\subsection{Specializations}\label{section4.3} Specialization morphisms with respect to perverse filtrations have been studied systematically in \cite{dCM, sp}. We provide a concrete description of the specialization morphism in our setting as follows.
\source{hitchin-fibrations-abelian-surfaces-pw}

As before, we assume $T = \BP^1$ and $T^\circ = \BP^1 \setminus 0$. Let $f: W \to T$ be a smooth morphism, whose restriction over $T^\circ$ is a trivial product
\[
f^\circ: W^\circ = W_t \times T^\circ \rightarrow T^\circ, \quad \forall~ t\neq 0,
\]
with $W_t$ proper.
By Deligne's Global Invariant Cycle Theorem, see  \cite[Theorem 1.7.1]{dCM1} for example, the restriction 
\[
\mathrm{res}_t:  H^*(W, \BQ) \to  H^*(W^\circ,\BQ)=H^*(W_t, \BQ)
\]
is surjective for $t \neq 0$. Let 
\[
\alpha_1, \alpha_2 \in H^*(W, \BQ)
\]
be two liftings of a class $\alpha \in H^*(W_t, \BQ)$ with $t\neq 0$. Since $H^*(W^\circ,\BQ)=H^*(W_t,\BQ)$, the long exact sequences for the pairs $(W,W_t)$ and $(W, W^\circ)$
are isomorphic. Since $W$ is nonsingular, the relative cohomology for $(W,W^\circ)$ can be identified with the Borel-Moore homology of $W_0$. It follows that
\begin{equation}\label{en52}
\source{hitchin-fibrations-abelian-surfaces-pw}
\alpha_1 - \alpha_2 = {i_0}_* \gamma, 
\end{equation}
for some class $\gamma$
 in the Borel--Moore homology of $W_0$.  Here $i_0: W_0 \hookrightarrow W$ is the closed embedding of the central fiber, and we identify the Borel--Moore homology and the cohomology of the nonsingular $W$ in the equation (\ref{en52}). The excess intersection formula \cite[Corollary 6.3]{Fulton}, {together with the triviality of the normal bundle $N_{W_0/W}$,} implies that we have
 \[
\mathrm{res}_0(\alpha_1) - \mathrm{res}_0(\alpha_2) = i_0^*{i_0}_* \gamma = c_1(N_{W_0/W})\cap \gamma =0.
\]
We define the \emph{specialization morphism}\footnote{In the literature, often, one defines a specialization morphism  $\mathrm{sp}^*$ in the opposite direction, as it is dictated by the morphism $\sigma: i_0^* \to \psi_f$ with $\psi_f$ nearby cycle-functor. In the present set-up, since $W_0$ is not proper, such an arrow $sp^*$ does not exist; 
by dualizing $\sigma$, one sees that an arrow in the opposite direction always exists, but if $W_0$ is singular, for example, this arrow is not a familiar object.}
\begin{equation}\label{sp!}
\source{hitchin-fibrations-abelian-surfaces-pw}
\mathrm{sp}^!: H^*(W_t, \BQ) \rightarrow H^*(W_0, \BQ), \quad t\neq 0
\end{equation}
as
\begin{equation}\label{sp1}
\source{hitchin-fibrations-abelian-surfaces-pw}
\mathrm{sp^!}(\alpha) = \mathrm{res}_0(\tilde{\alpha}) \in H^*(W_0, \BQ)
\end{equation}
where $\tilde{\alpha} \in H^*(W, \BQ)$ is any lifting of $\alpha$. The discussion above implies that the class (\ref{sp1}) does not depend on the lifting $\tilde{\alpha}$. Since $\mathrm{res}_0$ is a homomorphism of $\BQ$-algebras, so is $\mathrm{sp}^!$  (\ref{sp!}).


\begin{example}
\source{hitchin-fibrations-abelian-surfaces-pw}
\notready\label{example1}
\source{hitchin-fibrations-abelian-surfaces-pw}
Let $\CS \to T$ be the relative surface (\ref{normal_cone}) associated with the embedding $j_C: C \hookrightarrow A$. We have the specialization morphism
\begin{equation}\label{2333}
\source{hitchin-fibrations-abelian-surfaces-pw}
\mathrm{sp}^!: H^*(A, \BQ) \to H^*(T^*C, \BQ) = H^*(C, \BQ).
\end{equation}
By the definition of $\mathrm{sp}^!$, the morphism (\ref{2333}) is given by the pullback along
\[
T^*C  \hookrightarrow \BP(T^*C \oplus \CO_C) \hookrightarrow \mathrm{Bl}_{C\times 0}(A \times T) \to A\times T \to A.
\]
In particular, we have
\[
\mathrm{sp}^!(\gamma)  = j_C^* \gamma \in H^*(T^*C, \BQ) = H^*(C, \BQ),\quad \forall~ \gamma \in H^*(A, \BQ). 
\]
\end{example}

Now we discuss the interaction between $\mathrm{sp}^!$ and perverse filtrations. Let $W/T$ be as above. We consider a commutative diagram

\begin{equation*}
\begin{tikzcd}
W \ar [rr,"h"]\ar[rd]& & V\ar[ld]\\
 &T&
\end{tikzcd}
\end{equation*}  
where the morphism $h: W \to V$ is proper of relative dimension $d$. For every $t \in T$, there is a perverse filtration $P_\star H^*(W_t, \BQ)$ associated with the morphism $h_t : W_t \to V_t$.

\begin{prop}
\source{hitchin-fibrations-abelian-surfaces-pw}
\notready\label{Prop3.1}
\source{hitchin-fibrations-abelian-surfaces-pw}
The specialization morphism (\ref{sp!}) preserves\footnote{{In fact, the proof shows more, namely that $\mathrm{sp}^!$ is filtered strict.}} the perverse filtrations, \emph{i.e.}, 
\[
\mathrm{sp}^!\left( P_kH^*(W_t, \BQ ) \right) \subset P_k H^*(W_0, \BQ), \quad \forall k.
\]
\end{prop}

\begin{proof}
We show that there exist splittings of the perverse filtrations
\[
H^*(W_t, \BQ) = \bigoplus_{i}G_iH^*(W_t, \BQ), \quad \forall~ t \in T
\]
such that the splitting is constant for $t \in T^o$, and
\begin{equation*}
\mathrm{sp}^!\left(G_iH^*(W_t ,\BQ)\right) \subset G_iH^*(W_0, \BQ).
\end{equation*}




We apply the decomposition theorem \cite{BBD} to the morphism $h: W \to V$, and fix an isomorphism
\[
\phi: Rh_\ast \BQ_W [\dim{W} - d] \xrightarrow{\simeq} \bigoplus_{i=0}^{2d} \CP_V^i [-i].
\]
Here, the $\CP_V^i$ are semisimple perverse sheaves on $V$. By the discussion in Section \ref{sec1.3}, the isomorphism $\phi$ induces a splitting of the perverse filtration on $H^*(W, \BQ)$ associated with $h$. By the smoothness of $f$ and the compatibility of vanishing cycles with the derived pushforward $\mathrm{Rh}_*$, we have $\varphi(\mathrm{Rh}_\ast \BQ_W) = 0$.  As a result, \cite[Corollary 3.1.6]{dCM} shows that the restriction of $\CP_V^i$ to a closed fiber over $t\in T$ is a shifted perverse sheaf,
\[
\CP_{V,t}^i [-1] = \mathrm{res}_t^*\CP_V^i [-1] \in \mathrm{Perv}(V_t), \quad \forall~ t\in T.
\]
Therefore, the restriction $\phi_t$ of the decomposition $\phi$ induces  splittings of the perverse filtration
\[
H^*(W_t, \BQ) = \bigoplus_{i \geq 0} \phi_t (H^*(V_t, \CP_{V,t}^i)), \quad \forall~ t\in T,
\]
and we conclude from the definition of $\mathrm{sp}^!$ that
\[
\mathrm{sp}^! \left(\phi_t(H^k(V_t, \CP_{V,t}^i))\right) \subset \phi_0(H^k(V_0, \CP_{V,0}^i)).
\]
This completes the proof.
\end{proof}

Now we consider the family
\[
h_T: \CM \to  \CB
\]
over $T$ constructed in Section \ref{Section 4.2}, especially Lemma \ref{lem4.1}. Proposition \ref{Prop3.1} implies that the specialization morphism
\[
\mathrm{sp}^!: H^*(\CM_{\beta, A}, \BC) \to H^*(\CM_{\mathrm{Dol}}, \BC)
\]
preserves the perverse filtrations. Let $\widetilde{G}_*H^*(\CM_{A,\beta}, \BQ)$ be a splitting given by Theorem \ref{strengthened2.1}. In particular,  it is the first Deligne splitting induced by a  { $(\pi_\beta=h_{t\neq 0})$-}Lefschetz class\footnote{Note that this class was denoted by $\eta_{\beta',A'}$ in Section \ref{Section3.5}, since $A$ was special there.}
\[
\eta_{\beta,A} \in H^2(\CM_{A,\beta}, \BC).
\]



\begin{comment}
%associated with $h_t = \pi_\beta: \CM_{A,\beta} \to B$ and $h_0: \CM_{\mathrm{Dol}} \to \Lambda$.

%\begin{prop}
%Let $\hat{G}_*H^*(W_t, \BQ)$ be a splitting of the perverse filtration $P_*H^*(W_t, \BQ)$ for $t\neq 0$. Then we have the decomposition
%\[
%\mathrm{Im}\left\{\mathrm{sp}^!: H^*(W_t, \BQ) \to H^*(W_0,\BQ)\right\}= \bigoplus_{i,d} \mathrm{sp}^!\left(G_iH^d(W_0, \BQ)\right)
%\]
%which splits the restricted perverse filtration $P_*H^*(W_0, \BQ) \cap \mathrm{Im}(\mathrm{sp}^!)$.
%\end{cor}

\begin{proof}
Since $\mathrm{sp}^!$ preserves the perverse filtrations, it suffices to show that 
\[
\mathrm{sp}^!\left(G_iH^d(W_0, \BQ)\right) \cap \mathrm{sp}^!\left(G_jH^d(W_0, \BQ)\right) = \{0\}, \quad \forall~ i\neq j.
\]
This follows from Proposition \ref{Prop3.1} that
\[
\mathrm{sp}^!\left(G_iH^d(W_0, \BQ)\right) \cap \mathrm{sp}^!\left(P_{i-1}H^d(W_0, \BQ)\right) = \{0\}. \qedhere
\]
\end{proof}
\end{comment}

\begin{prop}
\source{hitchin-fibrations-abelian-surfaces-pw}
\notready\label{prop4.4}
\source{hitchin-fibrations-abelian-surfaces-pw}
The class
\[
\eta_0 = \mathrm{sp}^!\left( \eta_{\beta,A}\right) \in H^2(\CM_{\mathrm{Dol}}, \BC)
\]
is a Lefschetz class for the Hitchin fibration $h: \CM_{\mathrm{Dol}} \to \Lambda$. Assume further that \[
H^*(\CM_{\mathrm{Dol}}, \BC) = \bigoplus_{k,j}\widetilde{G}_k H^j(\CM_{\mathrm{Dol}}, \BC)\] is the first Deligne splitting associated with $\eta_0$. Then we have {(cf. Theorem \ref{strengthened2.1})}
\begin{equation}\label{prop4.4eq}
\source{hitchin-fibrations-abelian-surfaces-pw}
\mathrm{sp}^!\left( \widetilde{G}_kH^d(\CM_{A,\beta}, \BC) \right) \subset \widetilde{G}_k H^d(\CM_{\mathrm{Dol}}, \BC).\footnote{Here we use $\BC$-coefficients since the Lefschetz class $\eta_{A,\beta}$ lies in $H^2(\CM_{\beta,A}, \BC)$; see Section \ref{Section3.5}.}
\end{equation}
\end{prop}


\begin{proof}
Let $F_{\beta,A} \subset \CM_{\beta,A}$ be a 
closed fiber of the morphism $\pi_\beta: \CM_{\beta,A} \to B$. We have
\[
 \mathrm{dim}(F_{\beta,A})= \frac{1}{2}\mathrm{dim}(\CM_{\beta,A}) = \tilde{g}.
\]
Since the fiber class 
\[
[F_{\beta,A}] \in H^{2\tilde{g}}(\CM_{\beta,A}, \BC)
\]
lies in $P_0H^{2\tilde{g}}(\CM_{\beta,A}, \BC)$ by Proposition \ref{prop1.1}, and $\eta_{\beta,A}$ is a Lefschetz class, we obtain that
\[
\eta_{\beta,A}^{\tilde{g}}\cap{[F_{\beta,A}]} = \eta_{\beta,A}^{\tilde{g}} \cup [F_{\beta,A}] \neq 0
\]
by the hard Lefschetz condition. Hence after specialization, we have
\[
{\eta_0^{\mathrm{dim(F_0)}}} \cap [F_0]  = \eta_{\beta,A}^{\tilde{g}}\cap{[F_{\beta,A}]} \neq 0
\]
with $F_0$ a closed fiber of $h: \CM_{\mathrm{Dol}} \to \Lambda$. We conclude from Proposition \ref{prop3.2_Hitchin} that $\eta_0$ is a Lefschetz class with respect to $h$. The inclusion  (\ref{prop4.4eq}) then follows from Lemma \ref{comparison1} {(applied to $\phi = \mathrm{sp}^!$)}.
\end{proof}


\begin{rmk}
\source{hitchin-fibrations-abelian-surfaces-pw}
\notready\label{remark1}
\source{hitchin-fibrations-abelian-surfaces-pw}
Proposition \ref{prop4.4} implies that $\mathrm{sp}^!(G_*H^*(\CM_{\beta,A}, \BC))$ splits the restricted perverse filtration $P_*H^*(\CM_{\mathrm{Dol}}, \BC) \cap \mathrm{Im(sp^!)}$, \emph{i.e.},
\[
P_kH^*(\CM_{\mathrm{Dol}}, \BC) \cap \mathrm{Im(sp^!)} = \bigoplus_{i\leq k} \mathrm{sp}^!\left( \widetilde{G}_iH^*(\CM_{\beta,A}, \BC)\right).
\]
In general, for an arbitrary splitting $G_*H^*(\CM_{\beta,A}, \BC)$ of the perverse filtration $P_*H^*(\CM_{\beta,A}, \BC)$, it may not be true that
\[
\mathrm{sp}^!\left({G}_iH^*(\CM_{\beta,A}, \BC) \right) \cap \mathrm{sp}^!\left({G}_jH^*(\CM_{\beta,A}, \BC) \right) = \{0\}, \quad \forall i\neq j.
\]
It is thus crucial, in our approach to the results of this paper, to realize the splitting 
\[
H^*(\CM_{\beta,A}, \BC) = \bigoplus_{k,j}\widetilde{G}_kH^j(\CM_{\beta,A}, \BC)
\]
as the first  Deligne splitting associated with a Lefschetz class as in Theorem \ref{strengthened2.1}.
\end{rmk}



\subsection{Normalized classes} Our purpose is to apply the specialization morphism $\mathrm{sp}^!$ to the tautological classes. We first discuss some properties of the normalized classes introduced in Section \ref{sec_0.3}.

Following \cite{HRV, HT1, Shende}, the cohomology 
\begin{equation}\label{eqn89}
\source{hitchin-fibrations-abelian-surfaces-pw}
H^*(\CM_{\mathrm{Dol}}, \BQ) = H^*(\CM_B, \BQ)
\end{equation}
can be understood by the associated $\mathrm{PGL}_n$-character variety $\hat{\CM}_{\mathrm{B}}$. More precisely, let
\[
H^*(\CM_B, \BQ) \cong H^*(\hat{\CM}_{\mathrm{B}}, \BQ) \otimes H^*((\BC^*)^{2g}, \BQ)
\]
be the isomorphism established in \cite[Section 1]{HT1} and \cite[Theorem 2.2.12]{HRV}. Every class $w \in H^i(\hat{\CM}_{\mathrm{B}}, \BQ)$ can be naturally viewed as a class
\[
w = w\otimes 1 \in H^i(\CM_B, \BQ).
\]
The weights of the tautological classes associated with the universal $\mathrm{PGL}_r$-bundle $\CT$ on $C\times \CM_B$ (induced by the universal $\mathrm{GL}_r$-bundle) were calculated in \cite{Shende}, 
%{\MD J, do we mean $\hat{\CM}_{\mathrm{B}}$ below?}
\begin{equation}\label{eqn52}
\source{hitchin-fibrations-abelian-surfaces-pw}
\int_\gamma c_k(\CT) \in {^k\mathrm{Hdg}^{i+2k-2}(\CM_B)},\quad { \forall k \geq 0}, \forall~\gamma \in H^i(C, \BQ).
\end{equation}
The following lemma deduces (\ref{taut_weight}) from (\ref{eqn52}).

\begin{lem}
\source{hitchin-fibrations-abelian-surfaces-pw}
\notready\label{lem3.1}
\source{hitchin-fibrations-abelian-surfaces-pw}
A twisted universal family $(\CU^\alpha, \theta)$ on $C \times {\CM}_{\mathrm{Dol}}$ satisfies
\begin{equation}\label{53}
\source{hitchin-fibrations-abelian-surfaces-pw}
\int_\gamma \mathrm{ch}^\alpha_k(\CU) \in {^k\mathrm{Hdg}^{i+2k-2}(\CM_B)},\quad \forall~\gamma \in H^i(C, \BQ),~~ \forall~ k\geq 0,
\end{equation}
if and only if $\mathrm{ch}^\alpha(\CU)$ is normalized.
\end{lem}

\begin{proof}
We define
\[
{^k\mathrm{Hdg}^{d}(C\times \CM_B)} = \bigoplus_{i+j=d} H^i(C, \BQ) \otimes {^k\mathrm{Hdg}^{j}(\CM_B)}.
\]
Then (\ref{53}) is equivalent to
\[
\mathrm{ch}^\alpha(\CU) \in \bigoplus_k {^k\mathrm{Hdg}^{2k}(C\times \CM_B)}.
\]
By a direct calculation using Chern roots, we have (recall that $r$ is the rank)
\[
\mathrm{ch}^\alpha(\CU) = \mathrm{ch}(\CT) \cup \mathrm{exp}\left(\frac{c_1(\CU)}{r}+ \alpha\right).
\]
By using the universal relations between Chern character and total Chern class, we note that (\ref{eqn52}) is equivalent to \[
\mathrm{ch}(\CT) \in \bigoplus_k {^k\mathrm{Hdg}^{2k}(C\times \CM_B)}.
\]
By \cite{Shende}, we have
\[
H^0(\CM_B, \BQ) = {^0\mathrm{Hdg}^{0}(\CM_B)}, \quad H^2(\CM_B, \BQ) = {^2\mathrm{Hdg}^{2}(\CM_B)},
\]
which implies 
\[
{^1\mathrm{Hdg}^{2}(C\times \CM_B)} = H^1(C, \BQ)\otimes H^1(\CM_B ,\BQ).
\]
Therefore, we obtain that (\ref{53}) holds if and only if
\[
\mathrm{ch}_1^\alpha(\CU)= c_1(\CU) + r\alpha \in H^1(C, \BQ) \otimes H^1(\CM_B, \BQ).
\]
This is equivalent to the condition that the class $\mathrm{ch}(\CU^\alpha)$ is normalized.
\end{proof}

Next, in parallel to Lemma \ref{lem3.1},  we give a suffcient criterion involving the perverse filtration  for a class $\mathrm{ch}^\alpha(\CU)$ to be normalized.

\begin{lem}
\source{hitchin-fibrations-abelian-surfaces-pw}
\notready\label{lem4.7}
\source{hitchin-fibrations-abelian-surfaces-pw}
Suppose we have a twisted universal family $(\CU^\alpha, \theta)$ such that, for every $k \geq 0$ and 
every $\gamma \in H^{2i}(C, \BQ)$, the tautological class 
\begin{equation*}
\int_\gamma \mathrm{ch}^\alpha_k(\CU) \in H^{2i+2k-2}(\CM_\mathrm{Dol}, \BQ),
\end{equation*}
has perversity $k$. Then $\mathrm{ch}^\alpha(\CU)$ is normalized.
\end{lem}

\begin{proof}
We have 
\[
 P_1H^2(\CM_{\mathrm{Dol}}, \BQ) = 0, \quad  P_0H^0(\CM_{\mathrm{Dol}}, \BQ)=H^0(\CM_{\mathrm{Dol}}, \BQ).
\]
The first vanishing follows from the decomposition \eqref{112233}, and the corresponding vanishing for $\hat{\CM}_{\mathrm{Dol}}$ due to the fact that its second cohomology is generated by a relative ample class which must therefore have perversity $2$.\footnote{To see that a relative ample class has perversity 2, we apply \cite[Theorem 1.4.8]{dCHM1} together with the fact that an ample class does not vanish over a general fiber of the Hitchin morphism.} The second equality is clear in view of the isomorphism $H^0 (\CM_{\mathrm{Dol}},\BQ) \cong H^0 (\Lambda, \BQ)$ via pull-back from the base of the Hitchin fibration (\ref{Hitchin_fib}).
Since by assumption, any K\"unneth factor of $\mathrm{ch}_1^\alpha(\CU)$ in $H^*(\CM_{\mathrm{Dol}}, \BQ)$ has perversity $1$, we reach the desired conclusion
\[
\mathrm{ch}_1^\alpha(\CU) \in  H^1(C, \BQ)\otimes H^1(\CM_\mathrm{Dol} ,\BQ). \qedhere
\]
\end{proof}

\begin{rmk}
\source{hitchin-fibrations-abelian-surfaces-pw}
\notready\label{rmk4.8}
\source{hitchin-fibrations-abelian-surfaces-pw}
In fact, we see from the proofs of Lemmas \ref{lem3.1} and \ref{lem4.7} that both lemmas hold if we use $\BC$-coefficients for the cohomology groups and we allow $\alpha$ to be a $\BC$-class. In particular, a $\BC$-normalized class $\mathrm{ch}^\alpha(\CU)$ is unique and rational.
\end{rmk}



\subsection{Specializations and tautological classes}
Assume $j_C: C\hookrightarrow A$ is the closed embedding of the curve in an abelian surface. Let $\CS \to T$ and $\CM \to T$ be the family of surfaces and the relative moduli space of 1-dimensional sheaves introduced in Section \ref{Section 4.2}. In the following, we construct a family of cohomology classes on $\CM$ over $T$ whose restriction to every closed fiber $\CM_t$ is a twisted tautological class.

For our purpose, we take a relative compactification $\CS \subset \CS'$ over $T$ as in the proof of Lemma \ref{lem4.1} (ii). Then $\CM$ is the {coarse} moduli space of Gieseker-stable 1-dimensional sheaves $\CF$ on $\CS'$ satisfying that $\mathrm{supp}(\CF)$ is proper and contained in $\CS_t \subset \CS_t'$ with the numerical conditions
\[
\chi(\CF) = \chi, \quad [\mathrm{supp}(\CF)] = \beta_t.
\]
Here recall that $\beta_t$ is the class (\ref{beta_t}). By the proof of Lemma \ref{lem4.1} (ii), there is a universal $K$-theory class $[\CF_T]$ on $\CS' \times_T \CM$.

Let
\[
\pi_{\CS'}: \CS' \times_T \CM \to \CS', \quad \pi_\CM: \CS' \times_T \CM \to \CM
\]
be the projections. We have the Chern character 
\[
\mathrm{ch}({\CF_T}) \in H^*(\CS' \times_T \CM, \BC)
\]
defined by the universal class $[\CF_T]$. For any class of the type
\begin{equation}\label{alphatilde}
\source{hitchin-fibrations-abelian-surfaces-pw}
\tilde{\alpha} = \pi_\CS^* \tilde{\alpha}_1 + \pi_\CM^* \tilde{\alpha}_2 \in H^2(\CS' \times_T \CM, \BC)
\end{equation}
with $\tilde{\alpha}_1\in H^2(\CS', \BC),~ \tilde{\alpha}_2\in H^2(\CM, \BC)$, we consider the relative twisted class
\begin{equation}\label{456}
\source{hitchin-fibrations-abelian-surfaces-pw}
\mathrm{ch}^{\tilde{\alpha}}(\CF_T) = \mathrm{exp}(\tilde{\alpha})\cup \mathrm{ch}(\CF_T),
\end{equation}
whose degree $2k$ parts
\[
 \mathrm{ch}_k^{\tilde{\alpha}}(\CF_T) \in H^{2k}(\CS' \times_T \CM, \BC)
\]
induce a relative tautological class
\begin{equation}\label{427}
\source{hitchin-fibrations-abelian-surfaces-pw}
\int_{\tilde{\gamma}} \mathrm{ch}_k^{\tilde{\alpha}}(\CF_T): = {\pi_\CM}_* \left( (\pi_{\CS'}^*\tilde{\gamma} \cup  \mathrm{ch}_k^{\tilde{\alpha}}(\CF_T))\cap [\CS' \times_T \CM] \right) \in H_*^\mathrm{BM}(\CM,\BC) =H^*(\CM, \BC)
\end{equation}
for any $\tilde{\gamma}\in H^*(\CS', \BC)$. Here we use the Poincar\'e duality in the last identity, and we require the properness of $\pi_\CM$ for the pushforward functor ${\pi_\CM}_* : H^{\mathrm{BM}}_*(\CS' \times_T \CM, \BC) \to H^{\mathrm{BM}}_*(\CM, \BC)$.

Now we check that (\ref{427}) is the desired cohomology class on $\CM$ which gives a family of twisted tautological classes over the base $T$.


For any closed point $t \in T$, base change \cite[Proposition 1.7]{Fulton} implies that the restriction of (\ref{427}) to a closed fiber $\CM_t$ is a twisted universal class of the form 
\begin{equation}\label{restriction1}
\source{hitchin-fibrations-abelian-surfaces-pw}
\int_{\gamma_t} \mathrm{res}_t\left( \mathrm{ch}_k^{\tilde{\alpha}} \left(\CF_T\right)\right) \in H^*(\CM_t, \BC),\quad \gamma_t \in H^*(\CS'_t, \BC).
\end{equation}
When $t \neq 0$, we have $\CS_t= \CS'_t$, and therefore (\ref{restriction1}) recovers a twisted universal class on $\CM_{\beta,A}$,
\[
\int_{\gamma_t} \mathrm{ch}_k^\alpha (\CF_\beta) \in H^*(\CM_{\beta,A}, \BC).
\]
We now calculate the class (\ref{restriction1}) for $t =0$. By \cite[Theorem 4.6.5]{HL}, there exists a coherent sheaf $\CF_0$ supported on $T^*C \times \CM_{\mathrm{Dol}}$ which represents a universal class $[\CF_0]$ on $\CS'\times \CM_{\mathrm{Dol}}$. Moreover, the coherent sheaf $\CF_0$ is supported on the universal spectral curve $\CC \subset T^*C \times \CM_{\mathrm{Dol}}$ which is proper over $\CM_{\mathrm{Dol}}$. Hence by \cite[Property 2.1]{BFM}, we have
\begin{equation}\label{support123}
\source{hitchin-fibrations-abelian-surfaces-pw}
\mathrm{ch}_k(\CF_0) \cap [\CS_0'\times \CM_{\mathrm{Dol}}] = \sum_i c_i [Z_i] \in \mathrm{CH}_*(\CS'_0 \times \CM_{\mathrm{Dol}})_\BQ
\end{equation}
where $c_i \in \BQ$ and $Z_i$ are closed subvarieties in $T^*C \times \CM_{\mathrm{Dol}}$ which are proper over $\CM_{\mathrm{Dol}}$ via the composition
 \begin{equation}\label{composition1}
\source{hitchin-fibrations-abelian-surfaces-pw}
 Z_i \hookrightarrow T^*C \times \CM_{\mathrm{Dol}} \xrightarrow{p_\CM} \CM_{\mathrm{Dol}}. 
 \end{equation}
 
 
 
In particular, the expression (\ref{support123}) together with the properties for $Z_i$ implies that the class 
\[
\int_{\gamma_0} \mathrm{res}_0\left( \mathrm{ch}_k^{\tilde{\alpha}} \left(\CF_T\right)\right) \in H^*(\CM_{\mathrm{Dol}}, \BC)
\]
only depends on the restrictions of the classes {$\gamma_0$  to $H^2(T^*C, \BC)$ and
 $\mathrm{res}_0(\tilde{\alpha})$ to  $H^2(T^*C, \BC)\oplus H^2(\CM_{\mathrm{Dol}}, \BC)$}. 
 
 Note that a class of the type
\begin{equation}\label{820}
\source{hitchin-fibrations-abelian-surfaces-pw}
{p_\CM}_*(\omega \cup \mathrm{ch}_k(\CF_0)): = \sum_i c_i\cdot {p_{\CM}}_*(\omega \cap [Z_i]) \in H^{\mathrm{BM}}_*(\CM_{\mathrm{Dol}}, \BQ) = H^*(\CM_{\mathrm{Dol}}, \BQ)
\end{equation}
is well-defined for $\omega \in H^*(T^*C \times \CM_{\mathrm{Dol}}, \BC)$ due to the properness of (\ref{composition1}). We see from the discussion above that the restriction of the class (\ref{427}) to $\CM_0 (= \CM_{\mathrm{Dol}})$ is given by the following:
\begin{equation}\label{1114}
\source{hitchin-fibrations-abelian-surfaces-pw}
\mathrm{res}_0\left( \int_{\tilde{\gamma}} \mathrm{ch}_k^{\tilde{\alpha}} (\CF_T) \right) = \int_{\mathrm{sp}^!(\gamma_t)} \mathrm{ch}_k^{\alpha_0}(\CF_0) \in H^*(\CM_{\mathrm{Dol}}, \BC).
\end{equation}
Here $\mathrm{sp}^!(\gamma_t) \in H^*(\CS_0, \BC) = H^*(T^*C, \BC)$, $\alpha_0 \in H^2(T^*C, \BC)\oplus H^2(\CM_{\mathrm{Dol}}, \BC)$, and the class on the r.h.s. of (\ref{1114}) is defined by (\ref{820}).


The following proposition shows that the tautological classes (\ref{taut_class}) on the Dolbeault side $\CM_{\mathrm{Dol}}$  are obtained by specializing certain other tautological classes on the compact geometry side $\CM_{\beta,A}$.

\begin{prop}
\source{hitchin-fibrations-abelian-surfaces-pw}
\notready\label{prop4.8}
\source{hitchin-fibrations-abelian-surfaces-pw}
Let 
\[
\int_{\gamma} \mathrm{ch}_k^\alpha (\CF_\beta) \in H^*(\CM_{\beta,A}, \BC)
\]
be the classes of Theorem \ref{taut_abelian} with $\gamma \in H^*(A, \BQ)$ a rational class on $A$, then we have 
\[
\mathrm{sp}^!\left( \int_{\gamma} \mathrm{ch}_k^\alpha (\CF_\beta) \right) = c(k-1, j_C^*\gamma) \in H^*(\CM_{\mathrm{Dol}}, \BQ).
\]
\end{prop}

\begin{proof}
By the definition of $\mathrm{sp}^!$ and (\ref{1114}), we have
\[
\mathrm{sp}^!\left( \mathrm{res}_t\left( \int_{\tilde{\gamma}} \mathrm{ch}_k^{\tilde{\alpha}} (\CF_T) \right)\right) = \mathrm{res}_0\left( \int_{\tilde{\gamma}} \mathrm{ch}_k^{\tilde{\alpha}} (\CF_T) \right)= \int_{\mathrm{sp}^!(\gamma_t)} \mathrm{ch}_k^{\alpha_0}(\CF_0)
\]
for any $\tilde{\alpha}$ of the type (\ref{alphatilde}). 

%Hence the class $\mathrm{sp}^!\left( \int_{\gamma} \mathrm{ch}_k^\alpha (\CF_\beta) \right)$ is of the form {\MD see above $(***)$}
%\[
%\int_{\mathrm{sp}^!(\gamma)} \mathrm{ch}_k^{\alpha_0} (\CF_0),
%\]
%where $(\CF_0, \alpha_0)$ is a universal family of Gieseker stable 1-dimensional sheaves on $T^*C \times \CM_{\mathrm{Dol}}$.

A direct calculation by applying the Grothendieck--Riemann--Roch formula to the natural projection 
\[
\mathrm{pr}: T^*C\times \CM_{\mathrm{Dol}} \to C\times \CM_{\mathrm{Dol}} 
\]
(see the paragraph before \cite[Remark 8]{Markman3}) together with Example \ref{example1} yields
\begin{equation*}\label{997}
\source{hitchin-fibrations-abelian-surfaces-pw}
\int_{\mathrm{sp}^!(\gamma)} \mathrm{ch}_k^{\alpha_0} (\CF_0) = \int_{j_C^*\gamma} \mathrm{ch}_{k-1}^{\alpha_0}(\CU).
\end{equation*}
Here $\CU = \mathrm{pr}_* \CF_0$ is a universal family on $C\times \CM_{\mathrm{Dol}}$.

In conclusion, we obtain that
\[
\mathrm{sp}^!\left( \int_{\gamma} \mathrm{ch}_k^\alpha (\CF_\beta) \right) = \int_{j_C^*\gamma} \mathrm{ch}_{k-1}^{\alpha_0}(\CU).
\]
Moreover, we know from Proposition \ref{prop4.4} that the twisted class \[
\int_{j_C^*\gamma} \mathrm{ch}_{k-1}^{\alpha_0}(\CU),\quad \forall \gamma \in H^*(A, \BQ),~ \forall k\geq 1
\]
has perversity $k-1$. Hence Lemma \ref{lem4.7} and Remark \ref{rmk4.8} imply, by recalling (\ref{taut_class}), that 
\[
\int_{j_C^*\gamma} \mathrm{ch}_{k-1}^{\alpha_0}(\CU) = c(k-1, j_C^*\gamma). \qedhere
\]
\end{proof}

\subsection{Proofs of Theorems \ref{genus_2}, \ref{P=W_even}, and \ref{P=W_odd}}\label{Section4.6}
\source{hitchin-fibrations-abelian-surfaces-pw}

Since the perverse filtration with respect to the Hitchin fibration $h: \CM_{\mathrm{Dol}} \to \Lambda$ is locally constant when we deform the curve $C$ \cite{dCM}, it suffices to prove all the three theorems for a curve which can be embedded in an abelian surface
\[
j_C: C\hookrightarrow A.
\]
After having done so,  we can apply the specialization morphism 
\[
\mathrm{sp}^!: H^*(\CM_{\beta, A}, \BC) \to H^*(\CM_{\mathrm{Dol}}, \BC)
\]
introduced in Section \ref{section4.3}. The following results are immediate consequences of  Theorem \ref{strengthened2.1}, Proposition \ref{prop4.4}, and Proposition \ref{prop4.8}:
\begin{enumerate}
    \item[(i)] We have
    \begin{equation}\label{333333}
\source{hitchin-fibrations-abelian-surfaces-pw}
    c(\gamma, k) \in \widetilde{G}_kH^*(\CM_{\mathrm{Dol}}, \BQ), \quad \forall \gamma \in \mathrm{Im}({j_C^*}) \subset H^*(C, \BQ).
    \end{equation}
    \item[(ii)] The restriction of the decomposition \[
    H^*(\CM_{\mathrm{Dol}}, \BQ)= \bigoplus_{k,d}\widetilde{G}_kH^d(\CM_{\mathrm{Dol}}, \BQ)
    \]to the subalgebra of $H^*(\CM_{\mathrm{Dol}}, \BQ)$ generated by the classes (\ref{333333}) is multiplicative.
\end{enumerate}

We first prove Theorem \ref{genus_2}. The Abel-Jacobi morphism embeds a genus 2 curve $C$ into its Jacobian
\[
j_C: C \hookrightarrow \mathrm{Jac}(C) = A.
\]
Hence the restriction morphism $j_C^*$ is surjective, and Theorem \ref{genus_2} follows from (i) and (ii) above. 

The proof of Theorem \ref{P=W_even} is similar. For any embedding $j_C$, the image of $j_C^*$ always contains the sub-vector space
\[
H^0(C, \BQ) \oplus H^2(C, \BQ) \subset H^*(C, \BQ).
\]
Hence the subalgebra
\[
R^*(\CM_{\mathrm{Dol}}) \subset H^*(\CM_{\mathrm{Dol}}, \BQ)
\]
is contained in the subalgebra generated by the classes (\ref{333333}), and we again conclude Theorem \ref{P=W_even} by (i) and (ii).

Finally we treat the odd classes 
\[
c(\gamma, k) \in H^{2k-1}(\CM_{\mathrm{Dol}}, \BQ), \quad \gamma\in H^1(C, \BQ)
\]
and prove Theorem \ref{P=W_odd}. 

When the curve $C$ has genus $\geq 3$, the restriction
\begin{equation}\label{resH1}
\source{hitchin-fibrations-abelian-surfaces-pw}
j_C^*: H^1(A, \BQ) \to H^1(C, \BQ)
\end{equation}
is not surjective. We know from (i) that $c(\gamma, k)$ has perversity $k$ for any $\gamma$ lying in the image of (\ref{resH1}). Since the monodromy group of the moduli space $\CM_g$ of nonsingular genus g curves is the full symplectic group $\mathrm{Sp}_{2g}$ by \cite{A}, the sub-vector space 
\[
\mathrm{Im}\left({j_C^*}:H^1(A, \BQ)\to H^1(C,\BQ) \right) \subset H^1(C, \BQ)
\]
generates the total cohomology $H^1(C, \BQ)$ via the action of the monodromy group. We deduce Theorem \ref{P=W_odd} from \cite{dCM}. \qed



\begin{thebibliography}{10}

\bibitem{At} M. F. Atiyah, {\em K-theory,} Lecture notes by D. W. Anderson. W. A. Benjamin, Inc., New York-Amsterdam 1967.


\bibitem{A} N. A'Campo, {\em Tresses, monodromie et le groupe symplectique,} Comment. Math. Helv. 54 (1979), no. 2, 318--327.


\bibitem{BFM} P. Baum, W. Fulton and R. MacPherson, {\em Riemann-Roch for singular varieties,} Inst. Hautes Etudes Sci. Publ. Math. 45 (1975), 101--145.


\bibitem{B} A. Beauville, {\em Vari\'et\'es k\"ahl\'eriennes dont la premi\`ere classe de Chern est nulle,} J. Differential Geom. 18 (1983), no. 4, 755--782 (1984).

\bibitem{Be} A. Beauville, {\em Syst\`emes hamiltoniens compl\`etement int\'egrables associ\'es aux surfaces~$K3$,} Problems in the theory of surfaces and their classification (Cortona, 1988), 25--31, Sympos. Math., XXXII, Academic Press, London, 1991.

\bibitem{BNR} A. Beauville, M.S. Narasimhan, and S. Ramanan, {\em Spectral curves and the generalized theta divisor,}  J. Reine Angew. Math. (1989) 398, 169--179.


\bibitem{BBD} A. A. Be\u{\i}linson, J. Bernstein, and P. Deligne, {\em Faisceaux pervers,} Analysis and topology on singular spaces, I (Luminy, 1981), 5--171, Ast\'erisque, 100, Soc. Math. France, Paris, 1982.

\bibitem{Borel} A. Borel, {\em Density properties for certain subgroups of semi-simple groups without compact components,} Ann. of Math. (2) 72 (1960), 179--188.

%\bibitem{Bo} F. A. Bogomolov, {\em On the cohomology ring of a simple hyper-K\"ahler manifold (on the results of Verbitsky),} Geom. Funct. Anal. 6 (1996), no. 4, 612--618.

%\bibitem{BK} F. Bogomolov and N. Kurnosov, {\em Lagrangian fibrations for IHS fourfolds,} arXiv:\linebreak 1810.11011.

%\bibitem{B1} O. Biquard, {\em Fibr\'es de Higgs et connexions int\'egrables: le cas logarithmique (diviseur lisse),} Ann. Sci. \'Ecole Norm. Sup. (4) 30 (1997), no. 1, 41--96. 

%\bibitem{B2} O. Biquard, O. Garc\'ia-Prada, I. Mundet i Riera, {\em Parabolic Higgs bundles and representations of the fundamental group of a punctured surface into a real group,} arXiv:1510.04207.


\bibitem{CDP} W. Y. Chuang, D. E. Diaconescu, and G. Pan, {\em BPS states and the P=W conjecture,} Moduli spaces, 132--150, London Math. Soc. Lecture Note Ser., 411, Cambridge Univ. Press, Cambridge, 2014.

\bibitem{Park}  M. A. de Cataldo, {\em Perverse sheaves and the topology of algebraic varieties,} Geometry of moduli spaces and representation theory, 1--58, IAS/Park City Math. Ser., 24, Amer. Math. Soc., Providence, RI, 2017. 



\bibitem{dC} M. A. de Cataldo, {\em Hodge-theoretic splitting mechanisms for projective maps,} with an appendix containing a letter from P. Deligne, J. Singul. 7 (2013), 134--156.

\bibitem{sp} M. A. de Cataldo, {\em Perverse Leray filtration and specialization with applications to the Hitchin morphism,} preprint.


\bibitem{dCHM1} M. A. de Cataldo, T. Hausel, and L. Migliorini, {\em Topology of Hitchin systems and Hodge theory of character varieties: the case $A_1$,} Ann. of Math. (2) 175 (2012), no.~3, 1329--1407.

\bibitem{dCHM3} M. A. de Cataldo, T. Hausel, and L. Migliorini, {\em Exchange between perverse and weight filtration for the Hilbert schemes of points of two surfaces,} J. Singul. 7 (2013), 23--38.

\bibitem{dCM} M. A. de Cataldo and D. Maulik, {\em The perverse filtration for the Hitchin fibration is locally constant,} arXiv:1808.02235.

%\bibitem{dCM3} M. A. de Cataldo and L. Migliorini, {\em The Douady space of a complex surface,} Adv. in Math. 151 (2000), 283--312.

\bibitem{dCM3} M. A. de Cataldo and L. Migliorini, {\em The Chow groups and the motives of the Hilbert scheme of points on a surface,} Journal of Algebra 251, 824--848 (2002).

%\bibitem{dCM4} M. A. de Cataldo and L. Migliorini, {\em The Chow motive of semismall resolutions,} Mathematical Research Letters 11, 151--170 (2004).

\bibitem{dCM0} M. A. de Cataldo and L. Migliorini, {\em The Hodge theory of algebraic maps,} Ann. Sci. \'Ecole Norm. Sup. (4) 38 (2005), no. 5, 693--750.

\bibitem{dCM123} M. A. de Cataldo and L. Migliorini, {\em Intersection forms, topology of algebraic maps and motivic decompositions for resolutions of threefolds,} Algebraic cycles and motives, Vol. 1, 102--137, London Math. Soc. Lecture Note Ser., 343, Cambridge Univ. Press, Cambridge, 2007.

\bibitem{dCM1} M. A. de Cataldo and L. Migliorini, {\em The decomposition theorem, perverse sheaves and the topology of algebraic maps,} Bull. Amer. Math. Soc. (N.S.) 46 (2009), no. 4, 535--633.

\bibitem{dCM4} M. A. de Cataldo, L. Migliorini, {\em Hodge-theoretic aspects of the decomposition theorem,}
Algebraic Geometry, Seattle 2005, Proceedings of Symposia in Pure Mathematics, Vol.
80.2, 2009.



\bibitem{D} P. Deligne, {\em D\'ecompositions dans la cat\'egorie d\'eriv\'ee,} Motives (Seattle, WA, 1991), 115--128, Proc. Sympos. Pure Math., 55, Part 1, Amer. Math. Soc., Providence, RI, 1994.

\bibitem{DLL} R. Donagi, L. Ein, R. Lazarsfeld, {\em
Nilpotent cones and sheaves on K3 surfaces,} Birational algebraic geometry (Baltimore, MD, 1996), 51--61, Contemp. Math., 207, Amer. Math. Soc., Providence, RI, 1997. 

%\bibitem{D} P. Deligne, {\em Theorie de Hodge, III,} Publications Math\'ematiques de l'IHES 44.1 (1974): 5-77.

\bibitem{Fulton} W. Fulton, {\em Intersection theory,} Ergebnisse der Mathematik und ihrer Grenzgebiete (3), 2. Springer-Verlag, Berlin, 1984.



%\bibitem{Fuj} A. Fujiki, {\em On the de Rham cohomology group of a compact K\"ahler symplectic manifold,} Algebraic geometry, Sendai, 1985, 105--165, Adv. Stud. Pure Math., 10, North-Holland, Amsterdam, 1987.

%\bibitem{FH} W. Fulton and J. Harris, {\em Representation theory. A first course,} Graduate Texts in Mathematics, 129, Readings in Mathematics, Springer-Verlag, New York, 1991. xvi+551 pp.

%\bibitem{GV} R. Gopakumar and C. Vafa, {\em $M$-Theory and topological strings--II,} arXiv:hep-th/\linebreak 9812127.

%\bibitem{GNR} A. Gorsky, N. Nekrasov, V. Rubtsov, {\em Hilbert schemes, separated variables, and D-branes,} Comm. Math. Phys. 222 (2001), no. 2, 299--318.

\bibitem{Go1} L. G\"ottsche, {\em The Betti numbers of the Hilbert scheme of points on a smooth projective surface,} Math. Ann. 286 (1990), 193--207.

\bibitem{Go2} L. G\"ottsche, W. S\"orgel, {\em Perverse sheaves and the cohomology of Hilbert schemes of smooth algebraic surfaces,} Math. Ann. 296 (1993), 235--245.

%\bibitem{Gr} M. Gr\"ochenig, {\em Hilbert schemes as moduli of Higgs bundles and local systems,} Int. Math. Res. Not. (2014) 2014 (23): 6523--6575.

\bibitem{GLO} V. Golyshev, V. Luntz, and D. Orlov, {\em Mirror symmetry for abelian varieties,} J. Alg. Geom.10 (2001) 433--496.

\bibitem{HLSY} A. Harder, Z. Li, J. Shen, and Q. Yin, {\em P=W for Lagrangian fibrations and degenerations of hyper-K\"ahler manifolds,} Forum Math. Sigma., to appear.

\bibitem{HLR} T. Hausel, E. Letellier, and F. Rodriguez-Villegas, {\em Arithmetic harmonic analysis on character and quiver varieties,} Duke Math. J. 160 (2011), no. 2, 323--400.

\bibitem{HRV} T. Hausel and F. Rodriguez-Villegas, {\em Mixed Hodge polynomials of character varieties,} with an appendix by Nicholas M. Katz, Invent. Math. 174 (2008), no. 3, 555--624. 


%\bibitem{HT2} T. Hausel, M. Thaddeus, {\em Relations in the cohomology ring of the moduli space of rank 2 Higgs bundles,} J. Amer. Math. Soc. 16 (2003), no. 2, 303--327.

\bibitem{HT1} T. Hausel, M. Thaddeus, {\em Generators for the cohomology ring of the moduli space of rank 2 Higgs bundles,} Proc. London Math. Soc. (3) 88 (2004), no. 3, 632--658. 


\bibitem{HL} D. Huybrechts and M. Lehn, {\em The geometry of moduli spaces of sheaves,} Aspects of Mathematics, E31. Friedr. Vieweg and Sohn, Braunschweig, 1997. xiv+269 pp.


%\bibitem{Hit} N. J. Hitchin, {\em The self-duality equations on a Riemann surface,} Proc. London Math. Soc. (3) 55 (1987), no. 1, 59--126.

\bibitem{Hit1} N.J. Hitchin, {\em Stable bundles and integrable systems,} Duke Math. J. 54 (1987) 91--114.


%\bibitem{HST} S. Hosono, M.-H. Saito, and A. Takahashi, {\em Relative Lefschetz action and BPS state counting,} Internat. Math. Res. Notices 2001, no. 15, 783--816.

%\bibitem{HX} D. Huybrechts and C. Xu, {\em Lagrangian fibrations of hyperk\"ahler fourfolds,} arXiv:\linebreak 1902.10440.

%\bibitem{Hw} J.-M. Hwang, {\em Base manifolds for fibrations of projective irreducible symplectic manifolds,} Invent. Math. 174 (2008), no. 3, 625--644.

%\bibitem{I} M. Inaba, M. Saito, {\em Moduli of unramified irregular singular parabolic connections on a smooth projective curve,} Kyoto J. Math. 53, no. 2 (2013), 433--482.

%\bibitem{KKP} S. Katz, A. Klemm, and R. Pandharipande, {\em On the motivic stable pairs invariants of $K3$ surfaces,} with an appendix by R. P. Thomas, $K3$ surfaces and their moduli, 111--146, Progr. Math., 315, Birkh\"auser/Springer, Cham, 2016.

%\bibitem{KKV} S. Katz, A. Klemm, and C. Vafa, {\em $M$-theory, topological strings, and spinning black holes}, Adv. Theor. Math. Phys. 3 (1999), no. 5, 1445--1537.

%\bibitem{KY} T. Kawai and K. Yoshioka, {\em String partition functions and infinite products,} Adv. Theor. Math. Phys. 4 (2000), no. 2, 397--485. 

%\bibitem{KL} Y.-H. Kiem and J. Li, {\em Categorification of Donaldson--Thomas invariants via perverse sheaves,} arXiv:1212.6444v5.

%\bibitem{L} M. Lehn, {\em Chern classes of tautological sheaves on Hilbert schemes of points on surfaces,} Inventiones Mathematicae, 1999, 136(1):157--207.

%\bibitem{LS} M. Lehn, C. Sorger, {\em The cup product of the Hilbert scheme for K3 surfaces,} Inventiones mathematicae May 2003, Volume 152, Issue 2, pp 305--329.

%\bibitem{LQW} W. Li, Z. Qin, W. Wang, {\em Vertex algebras and the cohomology ring structure of Hilbert schemes of points on surfaces,} Math. Ann. 324 (2002), 105--133.

\bibitem{LL} E. Looijenga and V. A. Lunts, {\em A Lie algebra attached to a projective variety,} Invent. Math. 129 (1997), no. 2, 361--412. 

\bibitem{LP1} J. Le Potier, {\em Faisceaux semi-stables de dimension 1 sur le plan projectif,} Rev. Roumaine Math. Pures Appl., 38(7--8):635--678, 1993.

\bibitem{LP2} J. Le Potier, Syst\`emes coh\'erents et structures de niveau. Ast\'erisque No. 214 (1993).

\bibitem{Lieb} M. Lieblich, Moduli of sheaves: a modern primer. Algebraic geometry: Salt Lake City 2015, 361-388, Proc. Sympos. Pure Math., 97.2, Amer. Math. Soc., Providence, RI, 2018. 


\bibitem{Markman} E. Markman, {\em Generators of the cohomology ring of moduli spaces of sheaves on symplectic surfaces,} J. Reine Angew. Math. 544 (2002), 61--82. 


\bibitem{Markman5} E. Markman, {\em Integral generators for the cohomology ring of moduli spaces of sheaves over Poisson surfaces,}
Adv. Math., 208 (2007), 622-646.

\bibitem{Markman2} E. Markman, {\em On the monodromy of moduli spaces of sheaves on K3 surfaces,} J. Algebraic Geom.17(2008), 29--99.


\bibitem{Markman4} E. Markman, {\em Lagrangian fibrations of holomorphic-symplectic varieties of K3[n]-type,} In Anne Fr\"uhbis-Kr\"uger et. al., Algebraic and Complex Geometry, volume 71. Springer Proceedings in Math., 2014.


\bibitem{Markman3} E. Markman, {\em The monodromy of generalized Kummer varieties and algebraic cycles on their intermediate Jacobians,} arXiv:1805.11574.


%\bibitem{Ma1} D. Matsushita, {\em On fibre space structures of a projective irreducible symplectic manifold,} Topology 38 (1999), no. 1, 79--83. Addendum: Topology 40 (2001), no. 2, 431--432.

%\bibitem{Ma2} D. Matsushita, {\em Equidimensionality of Lagrangian fibrations on holomorphic symplectic manifolds,} Math. Res. Lett. 7 (2000), no. 4, 389--391.

%\bibitem{Ma3} D. Matsushita, {\em Higher direct images of dualizing sheaves of Lagrangian fibrations,} Amer. J. Math. 127 (2005), no. 2, 243--259.

%\bibitem{Ma4} D. Matsushita, {\em On deformations of Lagrangian fibrations,} $K3$ surfaces and their moduli, 237--243, Progr. Math., 315, Birkh\"auser/Springer, Cham, 2016.

%\bibitem{MPT} D. Maulik, R. Pandharipande, and R. Thomas, {\em Curves on $K3$ surfaces and modular forms,} with an appendix by A. Pixton, J. Topol. 3 (2010), no. 4, 937--996. 

\bibitem{MT} D. Maulik and Y. Toda, {\em Gopakumar--Vafa invariants via vanishing cycles,} Invent. Math. 213 (2018), no. 3, 1017--1097. 

\bibitem{Mellit} A. Mellit, {\em Cell decompositions of character varieties,} arXiv:1905.10685.


%\bibitem{Na} H. Nakajima, {\em Heisenberg algebra and Hilbert schemes of points on projective surfaces,} Ann. Math. 145 (1997), 379--388.

%\bibitem{NS} B. Nasatya, B. Steer, {\em Orbifold Riemann surfaces and the Yang-Mills-Higgs equations,} Annali della Scuola Normale Superiore di Pisa - Classe di Scienze 22.4 (1995): 595-643.

%\bibitem{Og} K. Oguiso, {\em Picard number of the generic fiber of an abelian fibered hyperk\"ahler manifold,} Math. Ann. 344 (2009), no. 4, 929--937.

%\bibitem{Ou} W. Ou, {\em Lagrangian fibrations on symplectic fourfolds,} J. Reine Angew. Math., to appear.

%\bibitem{PT} R. Pandharipande and R. P. Thomas, {\em 13/2 ways of counting curves,} Moduli spaces, 282--333, London Math. Soc. Lecture Note Ser., 411, Cambridge Univ. Press, Cambridge, 2014.

%\bibitem{PT2} R. Pandharipande and R. P. Thomas, {\em The Katz--Klemm--Vafa conjecture for $K3$ surfaces,} Forum Math. Pi 4 (2016), e4, 111 pp.

\bibitem{GIT} D. Mumford, Geometric Invariant Theory, Springer-Verlag, 1965.

\bibitem{SY} J. Shen and Q. Yin, {\em Topology of Lagrangian fibrations and Hodge theory and Hodge theory of hyper-K\"ahler manifolds,} appendix by Claire Voisin, Duke Math. J., to appear.

\bibitem{SZ} J. Shen and Z. Zhang, {\em Perverse filtrations, Hilbert schemes, and the P=W conjecture for parabolic Higgs bundles,} Algebr. Geom., to appear.

\bibitem{Shende} V. Shende, {\em The weights of the tautological classes of character varieties,} Int. Math. Res. Not. 2017, no. 22, 6832--6840.

%\bibitem{Simp88} C. T. Simpson, {\em Constructing variations of Hodge structure using Yang--Mills theory and applications to uniformization,} J. Amer. Math. Soc. 1 (1988), no. 4, 867--918. 

%\bibitem{Simp90} C. T. Simpson, {\em Harmonic bundles on noncompact curves,} J. Amer. Math. Soc. 3 (1990), no. 3, 713--770.


\bibitem{Simp} C. T. Simpson, {\em Higgs bundles and local systems,} Inst. Hautes \'Etudes Sci. Publ. Math. No. 75 (1992), 5--95.

\bibitem{Si1994II} C. T. Simpson, {\em Moduli of representations of the fundamental group of smooth projective varieties II,"} Inst. Hautes \'Etudes Sci. Publ. Math.   No. 80 (1994), 5-79.

%\bibitem{Simp95} C. T. Simpson, {\em The Hodge filtration on nonabelian cohomology,} Algebraic geometry---Santa Cruz 1995, 217--281, Proc. Sympos. Pure Math., 62, Part 2, Amer. Math. Soc., Providence, RI, 1997.

%\bibitem{Ver90} M. S. Verbitski\u{\i}, {\em Action of the Lie algebra of $SO(5)$ on the cohomology of a hyper-K\"ahler manifold,} Funct. Anal. Appl. 24 (1990), no. 3, 229--230.

%\bibitem{Ver95} M. Verbitsky, {\em Cohomology of compact hyperkaehler manifolds,} Thesis (Ph.D.)--Harvard University, 1995, 89 pp.

%\bibitem{Ver96} M. Verbitsky, {\em Cohomology of compact hyper-K\"ahler manifolds and its applications,} Geom. Funct. Anal. 6 (1996), no. 4, 601--611.

\bibitem{WSB} G. Williamson, {\em The Hodge theory of the decomposition theorem,} S\'eminaire Bourbaki Vol. 2015/2016, Expos\'es 1104--1119, Ast\'erisque No. 390 (2017), Exp. No. 1115, 335--367.

\bibitem{Yo} K. Yoshioka, {\em Moduli spaces of stable sheaves on abelian surfaces,} Math. Ann. 321 (2001),no. 4, 817--884.

\bibitem{Z} Z. Zhang, {\em Multiplicativity of perverse filtration for Hilbert schemes of fibered surfaces,} Adv. Math. 312 (2017), 636--679.

\end{thebibliography}
\end{document}




\begin{comment}
\begin{prop}
\source{hitchin-fibrations-abelian-surfaces-pw}
\notready
There exists a smooth family 
\[
\pi_\CM : \CM \to T
\]
together with a sheaf $\CF_\CM$ on $\CS \times_T \CM$ satisfying the following peoperties.
\begin{enumerate}
    \item[(a)] The restriction of $(\CS \times_T \CM, \CF_\CM)$ to the central fiber 
    \[
    \left( \pi_\CS \times_T \pi_\CM\right)^{-1}(0) \subset \CS \times_T \CM
    \]
    is $\left(T^*C \times \CM_{\mathrm{Dol}}, \CF_n \right)$ where $\CF_n$ is a universal family of 1-dimensional sheaves.
    \item[(b)] The restriction of $(\CS \times_T \CM, \CF_\CM)$ to 
    \[
    \left( \pi_\CS \times_T \pi_\CM\right)^{-1}(T^\circ) \subset \CS \times_T \CM
    \]
    is $\left((A\times \CM_{\beta,A}) \times T^\circ, \CF_\beta \boxtimes \CO_{T^\circ}\right)$.
\end{enumerate}
\end{prop}


We consider
\[
\CS \subset \bar{\CS} = \mathrm{Bl}_{C\times 0}(A \times \BP^1).
\]
The 3-fold $\bar{\CS}$ is nonsingular and projective. Let $\bar{\beta} \in H_2(\bar{\CS}, \BZ)$ be the strict transform of the curve class $n[C \times \mathrm{pt}]$ on $A \times \BP^1$. Assume that $\bar{\CM}$ is the moduli space of 1-dimensional Gieseker-stable sheaves $\CF$ (with respect to a polarization) on $\bar{\CS}$ satisfying that
\[
 [\mathrm{supp}(\CF)] = \bar{\beta}, \quad \chi(\CF) = 1.
\]
By \cite[Theorem 4.6.5]{HL}, there exists a universal family on $\bar{\CS} \times \bar{\CM}$, which is denoted by $\CF_{\bar{\CM}}$. Since $\bar{\beta}$ is a curve class supported on fibers of 
\begin{equation}\label{123}
\source{hitchin-fibrations-abelian-surfaces-pw}
\bar{\CS} = \mathrm{Bl}_{C\times 0}(A \times \BP^1)\rightarrow \BP^1,
\end{equation}
every $\CF \in \bar{\CM}$ is supported on a closed fiber of (\ref{123}) by the stability of $\CF$. Hence the universal family $\CF_{\bar{\CM}}$ is supported on the sub-scheme
\[
\bar{\CS} \times _{\BP^1} \bar{\CM} \subset \bar{\CS} \times \bar{\CM}.
\]

Now we define $\CM$ to be the sub-moduli space
\begin{equation}\label{CM}
\source{hitchin-fibrations-abelian-surfaces-pw}
\CM \subset \bar{\CM}
\end{equation}
parametrizing 1-dimensional sheaves in $\bar{\CM}$ which are supported on $\CS \subset \bar{\CS}$. The variety $\CM$ admits a morphism
\[
\pi_\CM: \CM \to T
\]
whose closed fiber over $t\in T$ is the moduli space of sheaves in $\CM$ which are supported in the closed fiber 
\[
\pi_\CS^{-1}(t) \subset \CS.
\]
The restriction of $\CF_{\bar{\CM}}$ gives a universal family $\CF'_\CM$ on $\CS \times_T \CM$, which satisfies the condition (a). Moreover, its restriction over $T^\circ \subset T$ is
\[
\CF'_\beta \boxtimes \CO_{T^\circ} \in \Coh\left((A\times  \CM_{\beta,A})\times T^\circ \right)
\]
for \emph{some} universal family $\CF'_{\beta}$ on $A\times \CM_{\beta,A}$. In particular, there exists a line bundle $\CL$ on $\CM$ such that
\[
\CF_\beta = \CF'_\beta \otimes p_\CM ^\ast \CL
\]
where $p_\CM: A \times \CM_{\beta,A} \to \CM_{\beta,A}$ is the projection. 

By spreading, we obtain a line bundle $\tilde{\CL}$ on $\CM$ whose restriction over $T^\circ$ is $\CL$. Let $\tilde{p}_{\CM}: \CS \times_T\CM \to \CM$ be the natural projection. Then the variety $\CM$ together and the sheaf 
\[
\CF_{\CM} = \CF'_\CM \otimes \tilde{p}_{\CM}^* \tilde{\CL}
\]
satisfies both the conditions (a) and (b). Not complete yet...
\end{comment}















\subsection{Fourier--Mukai transforms}\label{Sec2.5} Sections 2.5 to 2.7 are devoted to proving Theorem \ref{taut_abelian}. Let $E$ and $E'$ be two very general elliptic curves. We first treat the special case
\source{hitchin-fibrations-abelian-surfaces-pw}
\begin{equation}\label{special}
\source{hitchin-fibrations-abelian-surfaces-pw}
A= E \times E', \quad \beta = \sigma + nf
\end{equation}
where $\sigma = [\mathrm{pt} \times E']$ and $f = [E \times \mathrm{pt}]$ are the classes of a section and a fiber with respect to the elliptic fibration $p: A \to E'$.

By \cite[Theorem 3.15]{Yo}, we have an isomorphism of the moduli spaces
\[
\CM_{\beta,A} \cong M_{v_n, A}(= A^{[n]} \times \hat{A}),\quad v_n = (1,0,-n),
\]
induced by a relative Fourier--Mukai transform with respect to the elliptic fibration $p: A \to E'$. 

The support of a sheaf $\CF \in \CM_{\beta, A}$ is the union of a section and several fibers (with multiplicities). Hence the Hilbert--Chow morphism (\ref{Hilb-Ch}) is exactly the morphism
\begin{equation}\label{HC2}
\source{hitchin-fibrations-abelian-surfaces-pw}
p_n\times q: A^{[n]} \times \hat{A} \to E'^{(n)} \times E
\end{equation}
where we use $\hat{A}\cong A = E \times E'$ and $q$ is the projection to the first factor.  Under the identification
\[
A \times A = (E \times E) \times (E'\times E'),
\]
the kernel for the relative Fourier--Mukai transform associated with the elliptic fibration $p: A \to E'$ is given by
\begin{equation}\label{rel_P}
\source{hitchin-fibrations-abelian-surfaces-pw}
\CP'_E = \CP_E \boxtimes \CO_{\Delta_{E'}}.
\end{equation}
with
\[
\CP_E = \CO_{E\times E}(\Delta_E - \mathrm{pt}\times E - E\times \mathrm{pt})
\]
the normalized Poincar\'e line bundle. Recall the universal family $\CE_n$ on $A \times M_{v_n,A}$ defined in (\ref{univ_ideal}). This provides us a universal family on $A \times \CM_{\beta,A}$ of 1-dimensional Gieseker-stable sheaves,
\begin{equation}\label{univ2}
\source{hitchin-fibrations-abelian-surfaces-pw}
\CF_\beta = {R\pi_{13}}_\ast \left( \pi_{12}^\ast (\pi_1^*N\otimes{\CP'_{E}}) \otimes^{\BL} \pi_{23}^\ast \CE_n \right)
\end{equation}
where $\pi_i, \pi_{ij}$ are the projection from $A \times A \times M_{v_n,A}$ to the corresponding factors, and $N \in p^* \mathrm{Pic}(E')\subset \mathrm{Pic}(A)$.\footnote{Twisting by a line bundle $N$ pulled back from $E'$ along the morphism $p: A \to {E'}$ is to ensure that the (shifted) image of $\CO_A$ with respect to the relative Fourier--Mukai transform has Euler characteristic 1; see \cite[Section 3.2]{Yo}} We define the normalized universal family 
\[
\CF^{n}_\beta = {R\pi_{13}}_\ast \left( \pi_{12}^\ast {\CP'_{E}} \otimes^{\BL} \pi_{23}^\ast \CE_n \right) = \CF_\beta \otimes \pi_A^\ast N^{-1}.
\]

The universal family (\ref{univ2}) serves as the one in Theorem \ref{taut_abelian}. We construct in the following a decomposition
\begin{equation}\label{splitting1}
\source{hitchin-fibrations-abelian-surfaces-pw}
H^\ast(\CM_{\beta, A}, \BQ) = \bigoplus_{i,d} \widetilde{G}_iH^d(\CM_{\beta, A}, \BQ)
\end{equation}
which serves as the splitting in Theorem \ref{taut_abelian}.

A canonical splitting of the perverse filtration associated with
\[
q: \hat{A} =A \to E
\]
is given by
\[
G'_kH^d(\hat{A}, \BQ) = \bigoplus_{i+k =d} H^i(E, \BQ)\otimes H^k(E',\BQ).
\]
We define (\ref{splitting1}) as the decomposition
\begin{equation}\label{splitting2}
\source{hitchin-fibrations-abelian-surfaces-pw}
\widetilde{G}_kH^\ast(\CM_{\beta, A}, \BQ) = \bigoplus_{i+j=k}  \widetilde{G}_i H^\ast(A^{[n]}, \BQ) \otimes G'_jH^\ast(\hat{A},\BQ),
\end{equation}
where
$\widetilde{G}_\ast H^\ast(A^{[n]}, \BQ)$ was defined in Section \ref{Section1.5}.
Since $\widetilde{G}_\ast H^\ast(A^{[n]}, \BQ)$ splits the perverse filtration associated with 
\[
p_n:  A^{[n]} \to E'^{(n)},
\]
the filtration (\ref{splitting2}) splits the perverse filtration associated with the Hilbert--Chow morphism (\ref{Hilb-Ch}). 

The following theorem verifies Theorem \ref{taut_abelian} for the special case (\ref{special}), where we can choose the twisting class to be 
\[
\alpha = - \pi_A^* c_1(N) \in H^2(A\times \CM_{\beta,A}, \BQ).
\]

\begin{thm}
\source{hitchin-fibrations-abelian-surfaces-pw}
\notready\label{Step1}
\source{hitchin-fibrations-abelian-surfaces-pw}
The decomposition (\ref{splitting2}) is multiplicative and is compatible with the normalized tautological classes; that is, we have
\begin{equation}\label{1}
\source{hitchin-fibrations-abelian-surfaces-pw}
\int_\gamma \mathrm{ch}_k(\CF^n_\beta) \in \widetilde{G}_{k-1} H^\ast(\CM_{\beta,A}, \BQ),\quad \forall~\gamma \in H^\ast(A, \BQ).
\end{equation}
\end{thm}

\begin{proof}
The multiplicativity of (\ref{splitting2}) follows directly from Theorem \ref{taut_Hilb} (b). It suffices to prove (\ref{1}). 

We consider the multiplicative splitting 
\begin{equation}\label{eqn27}
\source{hitchin-fibrations-abelian-surfaces-pw}
\widetilde{G}_kH^\ast(A \times \CM_{\beta,A} ,\BQ) = H^\ast(A, \BQ) \otimes \widetilde{G}_kH^\ast(\CM_{\beta,A}, \BQ)
\end{equation}
on the total space $A \times \CM_{\beta,A}$. This decomposition splits the perverse filtration associated with 
\[
\mathrm{id}\times \pi_\beta: A \times \CM_{\beta, A} \to A \times B.
\]
The statement (\ref{1}) is equivalent to the following:
\begin{equation}\label{eqn28}
\source{hitchin-fibrations-abelian-surfaces-pw}
\mathrm{ch}(\CF^n_\beta) \in \bigoplus_{k\geq 1} \widetilde{G}_{k-1} H^{2k}(A\times \CM_{\beta, A}, \BQ).
\end{equation}

For a nonsingular projective variety $X$ with a decomposition $G_*H^*(X, \BQ)$ on its cohomology, we say that a class $\alpha \in H^*(X, \BQ)$ is \emph{balanced} with respect to this decomposition if 
\[
\alpha \in \bigoplus_{k} G_{k} H^{2k}(X, \BQ).
\]
We prove (\ref{eqn28}) by the following 3 steps.

\noindent\emph{Step 1.} We consider the new decomposition 
\begin{equation}\label{31}
\source{hitchin-fibrations-abelian-surfaces-pw}
\begin{split}
\overline{G}_k H^*(A \times \CM_{\beta, A}, \BQ)& = \bigoplus_{i+j=k} H^i(E, \BQ) \otimes H^*(E', \BQ) \otimes \widetilde{G}_jH^*(\CM_{\beta,A}, \BQ)\\
& = \bigoplus_{i+j=k} \widetilde{G}_i H^\ast(A \times A^{[n]}, \BQ) \otimes G'_jH^\ast(\hat{A}, \BQ).
\end{split}
\end{equation}
of the cohomology $H^\ast(A\times \CM_{\beta,A}, \BQ)$. It is a multiplicative decomposition splitting the perverse filtration associated with the morphism
\[
p \times \pi_\beta: A\times \CM_{\beta,A} \to E' \times B.
\]
We first show that the class
\[
\mathrm{ch}(\CE_n) \in H^\ast(A \times M_{A, v_n}, \BQ) = H^*(A \times \CM_{\beta,A}, \BQ)
\]
is balanced with respect to the splitting (\ref{31}).


Recall the expression (\ref{univ_ideal}) of the universal family $\CE_n$. By Theorem \ref{taut_Hilb} (a), the class
\[
\mathrm{ch}(\CI_n) \in  H^{\ast}(A\times A^{[n]}, \BQ)
\]
is balanced with respect to the splitting (\ref{splitting0}). Hence via pulling back along the projection
\[
\pi_{12}: A\times \CM_{\beta,A} = A\times A^{[n]} \times \hat{A} \to A \times A^{[n]},
\]
we obtain that the class 
\[
\begin{split}
\pi_{12}^\ast \mathrm{ch}(\CI_n) =  \mathrm{ch}(\CI_n) \otimes 1_{\hat{A}} & \in H^\ast(A \times A^{[n]}, \BQ) \otimes H^0(\hat{A}, \BQ)\\  & \subset H^*(A\times \CM_{\beta,A}, \BQ)
\end{split}
\]
is also balanced with respect to (\ref{31}). A direct calculation shows the class $\pi_{13}^\ast c_1(\CP)$ is balanced. Hence we conclude from the multiplicativity of (\ref{31}) that the class
\[
\mathrm{ch}(\CE_n) = \pi_{12}^\ast \mathrm{ch}(\CI_n) \otimes \pi_{13}^\ast \mathrm{ch}(\CP)=  \pi_{12}^\ast \mathrm{ch}(\CI_n) \otimes \pi_{13}^\ast \mathrm{exp}(c_1(\CP))
\]
is balanced. 


\noindent \emph{Step 2.} We further consider the multiplicative decomposition  
\begin{equation}\label{32}
\source{hitchin-fibrations-abelian-surfaces-pw}
\widetilde{G}_k H^*(A \times A \times \CM_{\beta, A}, \BQ) = H^*(A, \BQ) \otimes \overline{G}_k H^*(A \times \CM_{\beta, A}, \BQ)
\end{equation}
of $H^*(A \times A \times \CM_{\beta, A})$, which splits the perverse filtration associated with the morphism
\[
\mathrm{id} \times p \times \pi_\beta: A \times A \times \CM_{\beta, A} \to A \times E'\times B.
\]
We show that 
\begin{equation}\label{34}
\source{hitchin-fibrations-abelian-surfaces-pw}
\mathrm{ch}(\pi_{12}^\ast {\CP'_{E}} \otimes^{\BL} \pi_{23}^\ast \CE_n) \in \bigoplus_{k\geq 1}\widetilde{G}_{k-1} H^{2k}(A \times A\times \CM_{\beta, A}, \BQ)
\end{equation}

In fact, by Step 1, we obtain that the class
\[
\pi_{23}^* \mathrm{ch}(\CE_n) = 1_A \otimes \mathrm{ch}(\CE_n) \in H^0(A, \BQ) \otimes H^*(A\times \CM_{\beta,A}, \BQ)
\]
is balanced with respect to (\ref{32}). The relative Poincar\'e sheaf (\ref{rel_P}) can be expressed as
\begin{equation}\label{233}
\source{hitchin-fibrations-abelian-surfaces-pw}
\pi_{13}^\ast \mathrm{ch}(\CP'_E) = \pi_{13}^\ast \mathrm{ch}(\CP_E)\cup \pi_{13}^* \mathrm{ch}(\CO_{\Delta_{E'}}).
\end{equation}
The class $\pi_{12}^\ast \mathrm{ch}(\CP_E)$ is balanced via a direct calculation, and 
\[
\pi_{12}^*\mathrm{ch}(\CO_{\Delta_E'}) = \pi_{12}^* c_1(\CO_{\Delta_{E'}}) \in \widetilde{G}_0 H^2(A\times A \times \CM_{\beta,A}, \BQ).
\]
Hence by (\ref{233}) and the multiplicativity of (\ref{32}), we get
\[
\pi_{12}^* \mathrm{ch}(\CP'_E)\in \bigoplus_{k\geq 1} \widetilde{G}_{k-1}H^{2k}(A\times A \times \CM_{\beta,A}).
\]
This yields (\ref{34}) since
\[
\mathrm{ch}(\pi_{12}^\ast {\CP'_{E}} \otimes^{\BL} \pi_{23}^\ast \CE_n) = \pi_{12}^*\mathrm{ch}(\CP'_E) \cup \pi_{23}^\ast \mathrm{ch}(\CE_n).
\]

\noindent \emph{Step 3.} We finish the proof of (\ref{eqn28}). 

Appying the Grothendieck--Riemann-Roch formula to the projection
\[
\pi_{13}: A \times A \times \CM_{\beta, A} \to A \times \CM_{\beta,A},
\]
we obtain that
\[
\mathrm{ch}(\CF^n_\beta) = {\pi_{13}}_\ast \mathrm{ch}(\pi_{12}^\ast {\CP'_{E}} \otimes^{\BL} \pi_{23}^\ast \CE_n).
\]
Equivalently, the class $\mathrm{ch}(\CF^n_\beta)$ corresponds to the K\"unneth factor of the class
\[
\mathrm{ch}(\pi_{12}^\ast {\CP'_{E}} \otimes^{\BL} \pi_{23}^\ast \CE_n) 
\]
in the subspace
\[
H^4(A, \BQ) \otimes H^*(A \times \CM_{\beta,A}, \BQ) \subset H^*(A\times A \times \CM_{\beta, A}, \BQ).
\]
Here $H^4(A, \BQ)$ is spanned by the point class in the second factor of the product $A \times A \times \CM_{\beta, A}$. Hence (\ref{eqn28}) follows from (\ref{34}).
\end{proof}


\subsection{Perverse filtrations and $H^2(\CM_{\beta,A}, \BQ)$}
In this section, we assume that $A$ is any abelian surface with $\beta$ an ample curve class satisfying \[\beta^2= 2n \geq 4.\] Our goal is to prove a variant of Proposition \ref{prop1.1} for $\CM_{\beta,A}$.  More precisely, we introduce a \emph{canonical} class
 \begin{equation}\label{ample_base}
\source{hitchin-fibrations-abelian-surfaces-pw}
 L_\beta \in H^2(\CM_{\beta,A} ,\BQ)
 \end{equation}
to characterize the perverse filtration $P_{\ast} H^*(\CM_{\beta, A})$ associated with the morphism
\[
\pi_\beta: \CM_{\beta, A} \to B.
\]


Let $\CF_0 \in \CM_{\beta, A}$ be any given point. We consider the morphisms
\[
\mathrm{det}: \CM_{\beta,A} \to \hat{A}, \quad \CF \mapsto \mathrm{det}(\CF)\otimes \mathrm{det}(\CF_0)^\vee 
\]
and
\[
\hat{\mathrm{det}}: \CM_{\beta,A} \to A, \quad \CF \mapsto \mathrm{det}(\phi_\CP(\CF)) \otimes \mathrm{det}(\phi_\CP(\CF))^\vee.
\]
Here $\phi_\CP : D^b\Coh(A) \to D^b\Coh(\hat{A})$ is the Fourier--Mukai transform induced by the normalized Poincar\'e line bundle $\CP$ on $A \times \hat{A}$. By \cite{Yo}, the Albanese map for $\CM_{\beta,A}$ with respect to the reference point $\CF_0$ is
\[
\mathrm{alb} = \hat{\mathrm{det}} \times \mathrm{det}: \CM_{\beta,A} \to A \times \hat{A}.
\]
The closed fiber over the origin $0 \in A\times \hat{A}$, denoted by
\[
K_{\beta,A} = \mathrm{alb}^{-1}(0) \subset \CM_{\beta, A},
\]
is a holomorphic symplectic vairiety of Kummer type. It parametrizes 1-dimensional sheaves $\CF \in \CM_{\beta,A}$ satisfying
\[
 \mathrm{det}(\CF)= \mathrm{det}(\CF_0) \in \hat{A},\quad   \mathrm{det}(\phi_\CP(\CF)) = \mathrm{det}(\phi_\CP(\CF)) \in \hat{\hat{A}} = A.
\]
The restriction of (\ref{Hilb-Ch}) to the subvariety $K_{\beta, A}$ is a Lagrangian fibration
\[
\pi'_\beta: K_{\beta, A} \to  |\CO_A(\beta)| =  \BP^{n-1}.
\]
Here $ |\CO_A(\beta)|$ denotes the linear system associated with the divisor \[
\mathrm{supp}(\CF_0) \subset A.
\]

We now consider the following commutative diagram,
\begin{equation}\label{diagram1}
\source{hitchin-fibrations-abelian-surfaces-pw}
    \begin{tikzcd}
K_{\beta, A} \times A \times \hat{A}\arrow{r}{h} \arrow{d}{\pi'_\beta \times \mathrm{pr}_A} & \CM_{\beta,A}\arrow{d}{\pi_\beta} \\
 \BP^{n-1}\times A \arrow{r}{h'} & B.
\end{tikzcd} 
\end{equation}
Here the horizontal arrows are defined by
\[
h(\CF, a, L) = {t_a}_* \CF \otimes  L,\quad h'(C, a)= {t_a}_\ast C
\]
with $t_a: A \xrightarrow{\simeq} A$ the translation by $a$. Since both $h$ and $h'$ are finite and surjective, the pullback morphism
\begin{equation*}
    h^*: H^k(\CM_{\beta, A}, \BQ) \to H^k(K_{\beta, A} \times A \times \hat{A}, \BQ)
\end{equation*}
is injective. Moreover, it is an isomorphism when $k=2$ by G\"ottsche's formula \cite{Go1} for the Betti numbers of the generalized Kummer varieties.

When $n =2$, we have $\CM_{\beta,A} \cong A \times \hat{A}$ by \cite{Yo}, and therefore $K_{\beta,A}$ is a point. 

Now we consider the case $n \geq 3$. Recall the Mukai vector $v_\beta =(0,\beta ,1) \in S_A^+$. By \cite[Theorem 0.2 and (1.6)]{Yo}, the correspondence induced by the class 
\[
c_1(\CF_\beta) \in H^2(A \times \CM_{\beta,A})
\]
composed with the restriction map
\[
H^2(\CM_{\beta,A}, \BQ) \to H^2(K_{\beta,A}, \BQ)
\]
gives an isometry between $v_\beta^\perp \subset S_A^+\otimes \BQ$ and $H^2(K_{\beta,A}, \BQ)$,
\[
\theta: v_\beta^\perp  \xrightarrow{\simeq}  H^2(K_{\beta,A}, \BQ).
\]
Hence the pullback $h^*$ induces an isomorphism
\begin{equation}\label{eqn39}
\source{hitchin-fibrations-abelian-surfaces-pw}
\varphi: H^2(\CM_{\beta, A}, \BQ) \xrightarrow{\simeq} v_\beta^\perp \oplus H^2(A\otimes \hat{A}, \BQ),
\end{equation}
which does not depend on the choice of $\CF_0$ by the discussion of the second factor 
\[
q_{v_\beta}: H^2(\CM_{{\beta},A}, \BQ) \to H^2(A \times \hat{A}, \BQ)
\]
in \cite[Section 8, Page 40]{Markman3}.


We consider the following classes
\[
\begin{split}
    l_\beta & = \theta\left((0,0,1)\in v_\beta^\perp\right) \in H^2(K_{\beta,A} ,\BQ), \\
    L_\beta & = \varphi^{-1} \left( l_\beta \oplus (\beta\otimes 1_{\hat{A}}) \right) \in H^2(\CM_{\beta, A}, \BQ).
\end{split}
\]
The following proposition is the main result of this section.

\begin{prop}
\source{hitchin-fibrations-abelian-surfaces-pw}
\notready\label{prop2.4}
\source{hitchin-fibrations-abelian-surfaces-pw}
We have 
\[
P_kH^m(\CM_{\beta,A}, \BQ) = \left( \sum_{i\geq 1} \mathrm{Ker}(L_\beta^{n+k+i-m}) \cap \mathrm{Im}(L_\beta^{i-1}) \right) \cap H^m(\CM_{\beta, A}, \BQ).
\]
\end{prop}

Something to be added..









\begin{proof}
We denote by 
\[
\pi_K = \pi'_\beta \times \mathrm{pr}_A : K_{\beta, A} \times A \times \hat{A} \to \BP^{n-1}\times A
\]
the left vertical map of the diagram (\ref{diagram1}). 

By \cite[Example 3.1]{Markman4}, we have 
\[
l_\beta = {\pi'_\beta}^\ast \CO_{\BP^5}(1) \in H^2(K_{\beta,A}, \BQ).
\]
Hence 
\[
h^*L_\beta = l_\beta \oplus (\beta\otimes 1_{\hat{A}}) \in H^2(K_{\beta, A} \times A \times \hat{A},\BQ)
\]
is the pullback of an ample class on the base $\BP^{n-1}\times A$ via the morphism $\pi_K$. Since $h$ and $h'$ are finite, the $t$-exactness of finite morphisms implies that a class $\alpha \in H^*(\CM_{\beta, A}, \BQ)$ lies in $P_kH^*(\CM_{\beta}, A)$ if and only if
\[
h^* \alpha \in P_kH^*(K_{\beta, A} \times A \times \hat{A},\BQ), 
\]
where the perverse filtration for $K_{\beta, A} \times A \times \hat{A}$ is associated with the morphism $\pi_K$. Hence we obtain Proposition \ref{prop2.4} by applying Proposition \ref{prop1.1} to the morphism $\pi_K$ and the class $L = h^*L_\beta$.
\end{proof}

\subsection{Proof of Theorem \ref{taut_abelian}}
In this section, we complete the proof of Theorem \ref{taut_abelian} by combining the tools developed in the previous sections.

Let $A$ be an abelian surface and $\beta$ be an ample curve class. A \emph{perverse package} associated with the pair $(A, \beta)$ is defined to be a triple
\[
 (\CF_\beta, \alpha_\beta, L_\beta),
\]
where $\CF_\beta$ is a universal family on $A \times \CM_{\beta,A}$, $L_\beta \in H^2(\CM_{\beta,A}, \BC)$ is the class introduced in Section 2.6, and $\alpha \in H^2(A\times \CM_{\beta,A}, \BC)$ is of the type (\ref{alpha}).

An isomorphism between the perverse packages associated with $(A, \beta)$ and $(A', \beta')$ is the following:
\begin{enumerate}
    \item[(1)] A grading-preserving isomorphism of $\BC$-vector spaces 
    \[
    g: H^*(A, \BC) \xrightarrow{\simeq} H^*(A', \BC).
    \]
    \item[(2)] A graded isomorphism of $\BQ$-algebras
    \[
    \phi: H^*(\CM_{\beta,A}, \BC) \xrightarrow{\simeq} H^*(\CM_{\beta',A'}, \BC).
    \]
    \item[(3)] The isomorphisms $g$ and $\phi$ satisfy that
    \[
    \begin{split}
        (g\otimes \phi) \mathrm{ch}^{\alpha_\beta} (\CF_\beta) & = \mathrm{ch}^{\alpha_{\beta'}}(\CF_{\beta'}),\\
        \phi(L_\beta) & = L_{\beta'}.
    \end{split}
    \]
\end{enumerate}
Isomorphisms of perverse packages form an equivalence relation. We call two pairs $(A, \beta)$ and $(A', \beta')$ \emph{perversely isomorphic} if there is a perverse package associated with $(A,\beta)$ isomorphic to a perverse package associated with $(A', \beta')$.

\begin{prop}
\source{hitchin-fibrations-abelian-surfaces-pw}
\notready\label{prop2.5}
\source{hitchin-fibrations-abelian-surfaces-pw}
Assume that $(A, \beta)$ is perversely isomorphic to $(A', \beta')$. If Theorem \ref{taut_abelian} holds for $(A, \beta)$, then it also holds for $(A', \beta')$.
\end{prop}

\begin{proof}
By Proposition \ref{prop2.4}, the condition 
\[
\phi(L_\beta) = L_{\beta'}
\]
implies that the perverse filtrations on $H^*(\CM_{\beta,A}, \BC)$ and $H^*(\CM_{\beta', A'}, \BC)$ are identified via $\phi$. Let 
\[
H^*(\CM_{\beta,A}, \BC) = \bigoplus_{k,d} \widetilde{G}_k H^d(\CM_{\beta,A}, \BC)
\]
be the splitting of the perverse filtration for $\CM_{\beta, A}$ satisfying
\[
\int_\gamma \mathrm{ch}^{\alpha_\beta}_k(\CF_\beta) \in \widetilde{G}_{k-1}H^*(\CM_{\beta,A}, \BC).
\] 
Then the decomposition 
\[
\
\widetilde{G}_k H^d(\CM_{\beta',A'}, \BC) = \phi(\widetilde{G}_k H^d(\CM_{\beta,A}, \BC))
\]
also splits the perverse filtration for $\CM_{\beta',A'}$, and the corresponding tautological classes lie in its correct pieces by the condition (3).
\end{proof}

If $(A, \beta)$ is perversely isomorphic to $(A', \beta')$, then condition (2) implies that
\[
\beta^2 = \mathrm{dim}(\CM_{\beta,A}) -2 = \mathrm{dim}(\CM_{\beta',A'}) -2 = \beta'^2.
\]
The following theorem deduced from Theorem \ref{monodromy_main} verifies its converse.

\begin{thm}
\source{hitchin-fibrations-abelian-surfaces-pw}
\notready\label{thm2.6}
\source{hitchin-fibrations-abelian-surfaces-pw}
Any pairs $(A, \beta)$ and $(A', \beta')$ satisfying
\[
\beta^2 = \beta'^2 
\]
are perversely isomorphic.
\end{thm}

\begin{proof}

Our goal is to find a grading-preserving isomorphism of $\BC$-vector spaces
\begin{equation}\label{39}
\source{hitchin-fibrations-abelian-surfaces-pw}
g: H^*(A, \BC)  \to H^*(A', \BC)
\end{equation}
satisfying the following two properties:
\begin{enumerate}
    \item[(a)] $g$ preserve the intersection pairings,
    \item[(b)] the restriction of $g$ on $S_A^+\otimes \BQ$ satisfies
    \begin{equation}\label{40}
\source{hitchin-fibrations-abelian-surfaces-pw}
    \begin{split}
        g(0,0,1) & = (0,0,1),\\
        g(1,0,0) & = (1,0,0),\\
        g(0,\beta,0) & = (0, \beta',0).
    \end{split}
    \end{equation}
\end{enumerate}
and a graded ring isomorphism
 \[
 \phi : H^*(\CM_{\beta,A}, \BC) \to H^*(\CM_{\beta', A'}, \BC)
 \]
satisfying the condition (3).

We construct $(g, \phi)$ in two steps.

First, because the Mukai vectors $v_\beta$ and $v_{\beta'}$ are primitive of the same length,
we can apply \cite[Theorem 8.3]{Markman3} (which follows from Yoshioka \cite{Yo}) to find
\[ g_1: H^*(A, \BZ)  \to H^*(A', \BZ) ,\quad \phi_1: H^*(\CM_{\beta,A}, \BQ) \to H^*(\CM_{\beta', A'}, \BQ)\]
such that $g_1(v_\beta) = v_{\beta'}$ and
\begin{equation}\label{eq41}
\source{hitchin-fibrations-abelian-surfaces-pw}
(g_1 \otimes \phi_1) \mathrm{ch}(\CF_\beta) = \mathrm{ch}(\CF_{\beta'})
\end{equation}
with $\CF_{\beta'}$ a universal family on $A \times \CM_{\beta,A}$. 
If $\beta$ and $\beta'$ have different divisibility, then $g_1$ cannot satisfy condition b) above, so it requires further modification.

We now choose $g_2 \in \mathrm{SO}(S_{A'}^+\otimes \BC)_{v_{\beta'}}$ such that the restriction
$$g_2\circ g_1|_{S_{A}^{+}}: S_{A}^{+} \otimes \BC \rightarrow S_{A'}^{+}\otimes\BC$$
is grading-preserving and satisfies condition b).
Because we are working with complex coefficients, we can lift $g_2$ to an element of
$\mathrm{Spin}(S_{A'}^{+}\otimes \BC)_{v_{\beta'}}$ and take the associated isomorphisms
\[
\begin{split}
 g_2: &  H^*(A',\BC)  \rightarrow H^*(A',\BC)
\\
\phi_2=  \gamma_{g_2,v_{\beta'}}: &   H^*(\CM_{\beta',A'},\BC)  \rightarrow H^*(\CM_{\beta',A'},\BC). 
\end{split}
\]

We deduce from Theorem \ref{monodromy_main} that
\begin{equation}\label{eq42}
\source{hitchin-fibrations-abelian-surfaces-pw}
(g_2 \otimes \phi_2) \mathrm{ch}(\CF_{\beta'}) =  \mathrm{ch}^{\pi_\CM^*a}(\CF_{\beta'}), \quad a \in H^2(\CM_{\beta',A'}, \BC).
\end{equation}

Finally, we set $g = g_2 \circ g_1$ and $\phi = \phi_2\circ \phi_1$.
Then  (\ref{eq41}) and (\ref{eq42}) imply that
\[
(g\otimes \phi)\mathrm{ch}(\CF_{\beta}) = \mathrm{ch}^{\pi_\CM^*a}(\CF_{\beta'}).
\]
In particular, for any class $\alpha_\beta \in H^2(A\times \CM_{\beta,A}, \BC)$ of the type (\ref{alpha}), we have
\[
(g\otimes \phi)\mathrm{ch}^{\alpha_\beta}(\CF_{\beta}) = \mathrm{ch}^{\alpha_{\beta'}}(\CF_{\beta'}), \quad \alpha_{\beta'} = \pi_\CM^*a+ (g\otimes \phi)\alpha_\beta.
\]

Finally, under the identification (\ref{eqn39}), it follows from (\ref{40}) that the isomorphism $\phi$ induced by $g$ sends $(0,0,1) \in v_\beta^\perp$ to $(0,0,1) \in v_{\beta'}^\perp$ and sends $\beta\otimes 1_{\hat{A}} \in H^2(A\times \hat{A}, \BC)$ to $\beta' \otimes 1_{\hat{A'}} \in H^2(A'\times \hat{A'}, \BC)$. In particular, we have
\[
\phi(L_\beta) = L_{\beta'}.
\]
This completes the proof of Theorem \ref{thm2.6}. 
\end{proof}

Theorem \ref{taut_abelian} follows from Theorem \ref{Step1}, Proposition \ref{prop2.5}, and Theorem \ref{thm2.6}. Later in Section \ref{Section3.5}, we will prove a strengthened version of Theorem \ref{taut_abelian} which further requires the decomposition $\tilde{G}_*H^*(\CM_{\beta,A}, \BC)$ to be induced by a construction of Deligne. 

\begin{rmk}
\source{hitchin-fibrations-abelian-surfaces-pw}
\notready\label{remark2.8}
\source{hitchin-fibrations-abelian-surfaces-pw}
It was shown in \cite{dCM} that perverse filtrations behave well under deformations. However, for the pairs $(A, \beta)$ and $(A',\beta')$ such that $\beta$ and $\beta'$ have different divisibilities in $H^2(A, \BZ)$ and $H^2(A',\BZ)$, the fibrations $\pi_\beta$ and $\pi_{\beta'}$ are not deformation equivalent.

Hence, in order to reduce Theorem \ref{taut_abelian} for any pair $(A,\beta)$ to that for a special pair
\[
A= E \times E', \quad \beta = \sigma +n f,
\] 
it is essential to consider an equivalence relation weaker than deformation equivalences of fibrations $\pi_\beta$. This is our motivation to work with equivalences of perverse packages as introduced above. 
\end{rmk}
